\section{Residual dynamics}
\begin{definition}[Admissible descent trajectory]\label{def:admissible-descent-trajectory}
An admissible descent trajectory is a sequence of normalized residual states in which each successor lies in the admissible descent sector of its predecessor. If strict descent is chosen whenever available, the trajectory is called strictly selected.
\end{definition}

\begin{lemma}[First no-oscillation lemma]\label{lem:first-no-oscillation}
A strictly selected admissible descent trajectory cannot exhibit a two-step oscillation between distinct states if strict descent is available at both states.
\end{lemma}

\begin{proof}
Strict descent at the first state gives strict energy decrease; strict descent at the second state gives another strict decrease. Returning to the first state would therefore force the same energy to be both strictly larger and strictly smaller than itself, a contradiction.
\end{proof}

\begin{theorem}[Dynamic control theorem]\label{thm:dynamic-control}
Admissible descent dynamics on the normalized residual space has nonincreasing normalized weighted residual energy. Under strict descent, two-step oscillation is excluded; close-side behavior is represented by trapping close basins; promote-side behavior is represented by persistent honest promote regions; and every honest promote-to-close transition must pass through the residual frontier.
\end{theorem}

\begin{proof}
The monotonicity of energy follows from admissible descent. The no-oscillation claim is Lemma~\ref{lem:first-no-oscillation}. The basin and frontier statements are the dynamic reading of the previously defined geometric structures.
\end{proof}

